% 直角梯形 + 等腰梯形 对照图
\documentclass[tikz,border=4pt]{standalone}
\usepackage{ctex}
\usepackage{amsmath}
\usetikzlibrary{calc, decorations.markings}
\begin{document}
\begin{tikzpicture}[
  scale=0.85, thick,
  tickone/.style={postaction={decorate, decoration={markings,
    mark=at position 0.5 with {\draw (-1.5pt,-3pt) -- (1.5pt,3pt);}}}},
]
  % 直角梯形
  \begin{scope}[xshift=0cm]
    \coordinate (A) at (0,3);
    \coordinate (D) at (3,3);
    \coordinate (B) at (0,0);
    \coordinate (C) at (5,0);
    \draw (A)--(B)--(C)--(D)--cycle;
    % 两处直角
    \draw (B) ++(0.22,0) -- ++(0,0.22) -- ++(-0.22,0);
    \draw (A) ++(0.22,0) -- ++(0,-0.22) -- ++(-0.22,0);
    \node[above left]  at (A) {$A$};
    \node[above right] at (D) {$D$};
    \node[below left]  at (B) {$B$};
    \node[below right] at (C) {$C$};
    \node[below] at (1.5,-0.4) {\footnotesize 直角梯形};
  \end{scope}

  % 等腰梯形
  \begin{scope}[xshift=8cm]
    \coordinate (A) at (1,3);
    \coordinate (D) at (4,3);
    \coordinate (B) at (0,0);
    \coordinate (C) at (5,0);
    \draw[tickone] (A)--(B);
    \draw[tickone] (D)--(C);
    \draw (A)--(D);
    \draw (B)--(C);
    \node[above left]  at (A) {$A$};
    \node[above right] at (D) {$D$};
    \node[below left]  at (B) {$B$};
    \node[below right] at (C) {$C$};
    \node[below] at (2.5,-0.4) {\footnotesize 等腰梯形};
  \end{scope}
\end{tikzpicture}
\end{document}
